% GENERATED FILE. Edit canonical lessons, then run npm run generate:books.
% canonical-source-hash: fnv1a64:d05ed6e6ec84718f
% canonical-lessons: TE-C345-smasanam, TE-C345-samadhi, TE-C345-nivali, TE-C345-vardhanti, TE-C345-vilunama

\chapter{Remembering}
\label{ch:te-remembering}
\hlchaptermodality{345}

\begin{chapteropening}
\textbf{By the end of this chapter:} \emph{I can say a cremation ground, a grave, a tribute, a death anniversary and a will.}

Complete the last lesson of chapter 345: I can say a cremation ground, a grave, a tribute, a death anniversary and a will.
\end{chapteropening}

\section[\texorpdfstring{\te{శ్మశానం} (śmaśānaṁ)}{శ్మశానం (śmaśānaṁ)}]{\te{శ్మశానం} (śmaśānaṁ) — a cremation ground}

\label{lesson:TE-C345-smasanam}



\begin{quote}
Before the new one: say the Telugu for a funeral procession, then the Telugu for a funeral.
\end{quote}

\subsection*{You'll want to know: \te{శ్మశానం}}
\textbf{\te{శ్మశానం}} — \emph{śmaśānaṁ} — \textquotedblleft{}a cremation ground\textquotedblright{}.

\begin{quote}
\te{అంతిమయాత్ర} \te{శ్మశానం} \te{చేరింది}.
\end{quote}

\begin{quote}
\emph{antimayātra śmaśānaṁ cērindi.}
\end{quote}

\begin{quote}
The funeral procession reached the cremation ground.
\end{quote}

\begin{cousinweb}
\textbf{\te{శ్మశానం}} (\emph{śmaśānaṁ}) is borrowed from Sanskrit \emph{śmaśāna}, \textquotedblleft{}a place for burning the dead\textquotedblright{} (CDIAL 12658).
\end{cousinweb}

\begin{grammarlens}[title={What you've built}]
One more word for telling what happened.
\end{grammarlens}

\subsection*{Your turn}
\begin{itemize}
\raggedright
  \item \emph{Say it:} \emph{śmaśānaṁ}
  \item \emph{Say it:} \emph{śmaśānaṁ}, once more
  \item \emph{Say it:} \emph{antimayātra śmaśānaṁ cērindi.} — The funeral procession reached the cremation ground.
  \item \emph{Recall:} say the Telugu for a funeral procession, then the Telugu for a funeral, then say \emph{śmaśānaṁ} again
\end{itemize}

\begin{tcolorbox}[breakable,colback=teal!4,colframe=teal!35!black,title={Before you move on}]
What does \te{శ్మశానం} mean? (\textquotedblleft{}A cremation ground\textquotedblright{}.) Say it once more.
\end{tcolorbox}
\section[\texorpdfstring{\te{సమాధి} (samādhi)}{సమాధి (samādhi)}]{\te{సమాధి} (samādhi) — a grave; a tomb}

\label{lesson:TE-C345-samadhi}



\begin{quote}
Before the new one: say the Telugu for a funeral, then the Telugu for a cremation ground.
\end{quote}

\subsection*{You'll want to know: \te{సమాధి}}
\textbf{\te{సమాధి}} — \emph{samādhi} — \textquotedblleft{}a grave; a tomb\textquotedblright{}.

\begin{quote}
\te{ఆ} \te{పాత} \te{సమాధి} \te{దగ్గర} \te{ఒక} \te{చెట్టు} \te{ఉంది}.
\end{quote}

\begin{quote}
\emph{ā pāta samādhi daggara oka ceṭṭu undi.}
\end{quote}

\begin{quote}
There is a tree near that old tomb.
\end{quote}

Plural: \te{సమాధులు}, \emph{samādhulu}.

\begin{cousinweb}
\textbf{\te{సమాధి}} (\emph{samādhi}) is borrowed from Sanskrit \emph{samādhi} (Monier-Williams), whose first sense is \textquotedblleft{}putting together, joining\textquotedblright{}; in Telugu it names a tomb.
\end{cousinweb}

\begin{grammarlens}[title={What you've built}]
One more word for telling what happened.
\end{grammarlens}

\subsection*{Your turn}
\begin{itemize}
\raggedright
  \item \emph{Say it:} \emph{samādhi}
  \item \emph{Say it:} \emph{samādhi}, once more
  \item \emph{Say it:} \emph{ā pāta samādhi daggara oka ceṭṭu undi.} — There is a tree near that old tomb.
  \item \emph{Recall:} say the Telugu for a funeral, then the Telugu for a cremation ground, then say \emph{samādhi} again
\end{itemize}

\begin{tcolorbox}[breakable,colback=teal!4,colframe=teal!35!black,title={Before you move on}]
What does \te{సమాధి} mean? (\textquotedblleft{}A grave; a tomb\textquotedblright{}.) Say it once more.
\end{tcolorbox}
\section[\texorpdfstring{\te{నివాళి} (nivāḷi)}{నివాళి (nivāḷi)}]{\te{నివాళి} (nivāḷi) — a tribute}

\label{lesson:TE-C345-nivali}



\begin{quote}
Before the new one: say the Telugu for a cremation ground, then the Telugu for a grave.
\end{quote}

\subsection*{You'll want to know: \te{నివాళి}}
\textbf{\te{నివాళి}} — \emph{nivāḷi} — \textquotedblleft{}a tribute\textquotedblright{}.

\begin{quote}
\te{విద్యార్థులు} \te{తమ} \te{మాస్టారికి} \te{నివాళి} \te{అర్పించారు}.
\end{quote}

\begin{quote}
\emph{vidyārthulu tama māsṭāriki nivāḷi arpiñcāru.}
\end{quote}

\begin{quote}
The students paid tribute to their teacher.
\end{quote}

\te{నివాళి} \te{అర్పించు}, to pay tribute: the phrase every news report uses.

\begin{cousinweb}
\textbf{\te{నివాళి}} (\emph{nivāḷi}) is a Telugu word for a tribute, an offering of honour; its origin is not settled.
\end{cousinweb}

\begin{grammarlens}[title={What you've built}]
One more word for telling what happened.
\end{grammarlens}

\subsection*{Your turn}
\begin{itemize}
\raggedright
  \item \emph{Say it:} \emph{nivāḷi}
  \item \emph{Say it:} \emph{nivāḷi}, once more
  \item \emph{Say it:} \emph{vidyārthulu tama māsṭāriki nivāḷi arpiñcāru.} — The students paid tribute to their teacher.
  \item \emph{Recall:} say the Telugu for a cremation ground, then the Telugu for a grave, then say \emph{nivāḷi} again
\end{itemize}

\begin{tcolorbox}[breakable,colback=teal!4,colframe=teal!35!black,title={Before you move on}]
What does \te{నివాళి} mean? (\textquotedblleft{}A tribute\textquotedblright{}.) Say it once more.
\end{tcolorbox}
\section[\texorpdfstring{\te{వర్ధంతి} (vardhanti)}{వర్ధంతి (vardhanti)}]{\te{వర్ధంతి} (vardhanti) — a death anniversary}

\label{lesson:TE-C345-vardhanti}



\begin{quote}
Before the new one: say the Telugu for a grave, then the Telugu for a tribute.
\end{quote}

\subsection*{You'll want to know: \te{వర్ధంతి}}
\textbf{\te{వర్ధంతి}} — \emph{vardhanti} — \textquotedblleft{}a death anniversary\textquotedblright{}.

\begin{quote}
\te{నిన్న} \te{తాతగారి} \te{మొదటి} \te{వర్ధంతి} \te{జరిగింది}.
\end{quote}

\begin{quote}
\emph{ninna tātagāri modaṭi vardhanti jarigindi.}
\end{quote}

\begin{quote}
Yesterday was grandfather's first death anniversary.
\end{quote}

Families remember the dead on this day each year.

\begin{cousinweb}
\textbf{\te{వర్ధంతి}} (\emph{vardhanti}) comes from Sanskrit, but its exact source is not settled.
\end{cousinweb}

\begin{grammarlens}[title={What you've built}]
One more word for telling what happened.
\end{grammarlens}

\subsection*{Your turn}
\begin{itemize}
\raggedright
  \item \emph{Say it:} \emph{vardhanti}
  \item \emph{Say it:} \emph{vardhanti}, once more
  \item \emph{Say it:} \emph{ninna tātagāri modaṭi vardhanti jarigindi.} — Yesterday was grandfather's first death anniversary.
  \item \emph{Recall:} say the Telugu for a grave, then the Telugu for a tribute, then say \emph{vardhanti} again
\end{itemize}

\begin{tcolorbox}[breakable,colback=teal!4,colframe=teal!35!black,title={Before you move on}]
What does \te{వర్ధంతి} mean? (\textquotedblleft{}A death anniversary\textquotedblright{}.) Say it once more.
\end{tcolorbox}
\section[\texorpdfstring{\te{వీలునామా} (vīlunāmā)}{వీలునామా (vīlunāmā)}]{\te{వీలునామా} (vīlunāmā) — a will; a testament}

\label{lesson:TE-C345-vilunama}



\begin{quote}
Before the new one: say the Telugu for a tribute, then the Telugu for a death anniversary.
\end{quote}

\subsection*{You'll want to know: \te{వీలునామా}}
\textbf{\te{వీలునామా}} — \emph{vīlunāmā} — \textquotedblleft{}a will; a testament\textquotedblright{}.

\begin{quote}
\te{తాతగారు} \te{ఇల్లు} \te{నా} \te{పేరు} \te{మీద} \te{వీలునామా} \te{రాశారు}.
\end{quote}

\begin{quote}
\emph{tātagāru illu nā pēru mīda vīlunāmā rāśāru.}
\end{quote}

\begin{quote}
Grandfather wrote a will leaving the house to me.
\end{quote}

\te{వీలునామా} \te{రాయు}, to make a will.

\begin{cousinweb}
\textbf{\te{వీలునామా}} (\emph{vīlunāmā}) ends in Persian \emph{nāma}, \textquotedblleft{}a writing, a deed\textquotedblright{}, as several loans for documents do; the origin of the first part is not settled.
\end{cousinweb}

\begin{grammarlens}[title={What you've built}]
One more word for telling what happened.
\end{grammarlens}

\subsection*{Your turn}
\begin{itemize}
\raggedright
  \item \emph{Say it:} \emph{vīlunāmā}
  \item \emph{Say it:} \emph{vīlunāmā}, once more
  \item \emph{Say it:} \emph{tātagāru illu nā pēru mīda vīlunāmā rāśāru.} — Grandfather wrote a will leaving the house to me.
  \item \emph{Recall:} say the Telugu for a tribute, then the Telugu for a death anniversary, then say \emph{vīlunāmā} again
\end{itemize}

\begin{tcolorbox}[breakable,colback=teal!4,colframe=teal!35!black,title={Before you move on}]
What does \te{వీలునామా} mean? (\textquotedblleft{}A will; a testament\textquotedblright{}.) Say it once more.
\end{tcolorbox}
